\clearpage
\begin{summarybox}
\textbf{\textcolor{CSummary}{先认原有对称，再看新增对称}}

偶函数原有一条轴（$y$ 轴）；奇函数原有一个中心（原点）。
在 $a\ne0$、全域对称和平移条件成立时：再增加同类对称，移动 $2|a|$ 直接重复；再增加不同类对称，移动 $2|a|$ 先变号，再移动同样距离得到 $4|a|$ 的周期。
这里“同类、不同类”只比较原有的轴/中心与新增的轴/零高度中心，不能脱离条件机械套用。

\medskip
\textbf{\textcolor{CSummary}{不同类搭配：两翻一变号，再翻回原样}}

\textbf{奇函数 $+$ 关于 $x=a$ 对称}，或\textbf{偶函数 $+$ 关于 $(a,0)$ 对称}，其中 $a\ne0$，且函数定义在 $\mathbb{R}$ 上时：
\[
\boxed{f(x+2a)=-f(x)}\quad\Longrightarrow\quad\boxed{f(x+4a)=f(x)}.
\]
所以\textbf{$4|a|$ 是一个正周期，不保证是最小正周期}。

\medskip
\textbf{\textcolor{CSummary}{两步速答}}
\begin{enumerate}[leftmargin=*,itemsep=0.2em]
 \item 先认组合：奇配轴，偶配点 $\Rightarrow$ 走 $2a$ 变号。
 \item 再走一次：走 $4a$ 恢复 $\Rightarrow$ $4|a|$ 是周期；若求函数值，直接把大数按这个周期化小。
\end{enumerate}

\medskip
\textbf{\textcolor{CSummary}{四种搭配对照}}
\begin{center}
\renewcommand{\arraystretch}{1.4}
\begin{tabular}{c|c|c}
 & 新增轴 $x=a$ & 新增中心 $(a,0)$\\\hline
偶函数 & $f(x+2a)=f(x)$ & $f(x+2a)=-f(x)$\\
一个正周期 & $2|a|$ & $4|a|$\\\hline
奇函数 & $f(x+2a)=-f(x)$ & $f(x+2a)=f(x)$\\
一个正周期 & $4|a|$ & $2|a|$
\end{tabular}
\end{center}

\textbf{\textcolor{CSummary}{最容易错的一句话}}

“$4|a|$ 是周期”\textbf{不等于}“$4|a|$ 是最小正周期”。只给两条对称信息，不能无依据地判断最小正周期。$2|a|$ 同样只保证是一个正周期。$a<0$ 要取绝对值；$a=0$ 不能套出正周期；定义域须允许所有相关对称和来回平移，且等式对全域成立。

\medskip
\textbf{\textcolor{CSummary}{再多记一条好用的}}

奇配轴可得到 $f(0)=f(2a)=0$；偶配点可得到 $f(a)=f(3a)=0$。遇到求函数值或判断零点的题，往往比先求周期更快。
\end{summarybox}
